% ==============================================================================
% CHƯƠNG 2: PHÂN TÍCH YÊU CẦU VÀ ĐẶC TẢ HỆ THỐNG
% ==============================================================================

\chapter{Phân tích yêu cầu và Đặc tả hệ thống}
\label{chap:phan-tich-yeu-cau}

\section{Phân tích bài toán và nhu cầu người dùng}
\label{sec:phan-tich-bai-toan}

% Nội dung phần 2.1

\section{Phân tích tác nhân và Use Case}
\label{sec:phan-tich-tac-nhan}

% Nội dung phần 2.2

\section{Đặc tả yêu cầu chức năng}
\label{sec:dac-ta-chuc-nang}

\subsection{Quản lý người dùng và xác thực}
\label{subsec:quan-ly-nguoi-dung}

% Nội dung 2.3.1

\subsection{Tìm kiếm và khám phá research paper}
\label{subsec:tim-kiem-paper}

% Nội dung 2.3.2

\subsection{Quản lý và khai thác thông tin research paper}
\label{subsec:khai-thac-paper}

% Nội dung 2.3.3

\subsection{Tương tác và hỗ trợ tìm hiểu research paper bằng AI}
\label{subsec:tuong-tac-ai}

% Nội dung 2.3.4

\subsection{Quản lý lịch sử và dữ liệu người dùng}
\label{subsec:quan-ly-lich-su}

% Nội dung 2.3.5

\section{Đặc tả yêu cầu phi chức năng}
\label{sec:dac-ta-phi-chuc-nang}

\subsection{Hiệu năng và khả năng đáp ứng}
\label{subsec:hieu-nang}

% Nội dung 2.4: Hiệu năng và khả năng đáp ứng

\subsection{Bảo mật và quyền riêng tư}
\label{subsec:bao-mat}

% Nội dung 2.4: Bảo mật và quyền riêng tư

\subsection{Tính ổn định và độ tin cậy}
\label{subsec:do-tin-cay}

% Nội dung 2.4: Tính ổn định và độ tin cậy

\subsection{Khả năng bảo trì và mở rộng}
\label{subsec:kha-nang-bao-tri}

% Nội dung 2.4: Khả năng bảo trì và mở rộng

\section{Phân tích luồng nghiệp vụ và luồng xử lý}
\label{sec:luong-nghiep-vu}

% Nội dung phần 2.5

\section{Phân tích luồng dữ liệu}
\label{sec:luong-du-lieu}

% Nội dung phần 2.6

\section{Xác định tình huống lỗi và cơ chế xử lý}
\label{sec:tinh-huong-loi}

% Nội dung phần 2.7

\section{Xác định quy mô dữ liệu và khả năng mở rộng hệ thống}
\label{sec:quy-mo-du-lieu}

% Nội dung phần 2.8
